% =====================================================================
%  contact.tex — the Contact page
%  (the links on the left come from the \Social lines in site.tex)
% =====================================================================

\PageTitle[yellow]{Contact}

\begin{Intro}
Feel free to say hi. I usually reply within a few days. \texttt{:) }
\end{Intro}

\FormTitle{Contact Me}

% \FormField[type]{name}{Label}{Placeholder}     type: email, text, textarea
\FormField[email]{email}{Your mail:}{example@email.com}
\FormField[text]{subject}{Subject:}{About...}
\FormField[textarea]{message}{Content:}{Type here...}

\FormButton{Send}

% Without \FormEndpoint the Send button opens the visitor's mail app with the message filled in
% (to the mail address in site.tex); \FormNote is shown next to the button after that.
\FormNote{Your mail app should open with the message ready to send. If nothing happens, please write to the e-mail address on the left.}

% The message is sent straight from the page through Web3Forms (access key from web3forms.com;
% it is meant to be public and only allows sending mail to your own address).
% To go back to the "open the mail app" behaviour, put a % in front of the next two lines.
\FormEndpoint{https://api.web3forms.com/submit}
\FormAccessKey{cfbf2d1b-9443-4a54-b9cb-e968f21aa858}
\FormSent{Thank you, your message has been sent.}
\FormError{Sorry, the message could not be sent. Please write to the e-mail address on the left.}
